\documentclass[../../main.tex]{subfiles}
\begin{document}
\rhead{東京大学後期$2001$年$1$番}
\setcounter{equation}{0}
\begin{oframed}
  \subfile{problem}
\end{oframed}
\setlength{\columnseprule}{0.4pt}
\begin{multicols}{2}
  {\bf [解答]}
  表現の簡潔さのため
  \begin{align}
    & f(n) = \sum_{k=2}^{n}\dfrac{k}{\sqrt{k^2-1}}         \\
    & g(x) = \dfrac{x}{\sqrt{x^2-1}}               & (x>1)
  \end{align}
  とおくと，
  \begin{align}
    f(n) = \sum_{k=2}^{n} g(k) \label{eq:1}
  \end{align}
  である．さらに
  \begin{align}
    A(n) = n-f(n)
  \end{align}
  とおく．
  $A(n)$の評価をすることで，題意の不等式が成り立つ$i$を求める．

  まず$y=g(x)$のグラフの概形を調べる．
  一階，二階微分は
  \begin{align}
    g'(x)
    & = -1(x^2-1)^{-3/2} < 0 \\
    g''(x)
    & =3x(x^2-1)^{-5/2} > 0
  \end{align}
  だから，$g(x)$は下に凸で，
  \begin{align}
    & \lim_{x\to 1} g(x) \to \infty \\
    & \lim_{x\to \infty} g(x) \to 1
  \end{align}
  だから，$y=g(x)$のグラフは\cref{2001-1:fig:1}のようになる．

  \begin{figure}[H]
    \centering
    \begin{tikzpicture}
      \begin{axis}[
          xmin=0, xmax=12,
          ymin=0, ymax=3,
          axis lines=middle,
          xlabel=$x$,
          ylabel=$y$,
          xtick={1},
          ytick={1},
          grid=none,
          domain=1.2:10,
          samples=200,
          no markers,
          clip=false,
        ]

        % Draw the function y=x/sqrt(x^2-1)
        \addplot[thick, smooth, blue] {x/sqrt(x^2-1)} node[above, pos=0.5, yshift=5mm, black] {$y=\frac{x}{\sqrt{x^2-1}}$};

        % Draw the vertical line at x=1
        \draw[dashed, red] (axis cs:1, 0) -- (axis cs:1, 3);

        % Draw the horizontal line at y=1
        \draw[dashed, red] (axis cs:0, 1) -- (axis cs:5, 1);

        % Add a label for the origin
        \node[below left] at (axis cs:0, 0) {$O$};
      \end{axis}
    \end{tikzpicture}
    \caption{$y=g(x)$のグラフの概形．}
    \label{2001-1:fig:1}
  \end{figure}

  面積を比較するため，
  \begin{align}
    S_n = \int_{3}^{n}g(x) dx \label{eq:2}
  \end{align}
  とおくと，\cref{eq:1}に注意して$S_n,f(n)$の示す面積は\cref{2001-1:fig:2}の斜線部である．
  以下$n \ge 3$として考える．

  \pgfmathdeclarefunction{gfunc}{1}{%
    \pgfmathparse{#1/sqrt((#1)^2-1)}%
  }

  \begin{figure}[H]
    \centering
    \begin{tikzpicture}
      \begin{axis}[
          xmin=0, xmax=60,
          ymin=0, ymax=2.0,
          axis lines=middle,
          axis line style={->, thick, gray!80!black},
          xlabel=$x$,
          ylabel=$y$,
          xtick={1,3,50},
          xticklabels={1,3,$n$},
          ytick={1},
          grid=none,
          samples=200,
          no markers,
          clip=false,
        ]

        % Draw the function y=x/sqrt(x^2-1)
        \addplot[thick, smooth, blue!85!black,domain=1.16:60,samples=200] {gfunc(x)} node[above, pos=0.5, yshift=5mm, black] {$y=\frac{x}{\sqrt{x^2-1}}$};
        \addplot[draw=none, domain=3:50, samples=100, name path=A] {gfunc(x)} ;
        \addplot[draw=none, domain=3:50, name path=line]{0};
        % Draw the vertical line at x=1
        \draw[densely dotted, thick, gray!70!black] (axis cs:1, 0) -- (axis cs:1, 2);
        % Draw the horizontal line at y=1
        \draw[densely dotted, thick, gray!70!black] (axis cs:0, 1) -- (axis cs:50, 1);
        % Add a label for the origin
        \node[below left] at (axis cs:0, 0) {$O$};
        \draw[densely dotted, thick, gray!70!black] (axis cs:50, 0) -- (axis cs:50, 1);

        \addplot[pattern=north east lines, pattern color=gray!60!black] fill between [of=A and line];
      \end{axis}
    \end{tikzpicture}
    \caption{$S_n$の面積を図示したもの．}
    \label{2001-1:fig:2}
  \end{figure}

  \begin{figure}[H]
    \centering
    \pgfmathsetmacro\maxx{12}
    \pgfmathsetmacro\maxn{10}
    \begin{tikzpicture}
      \begin{axis}[
          xmin=0, xmax=\maxx,
          ymin=0, ymax=2.0,
          axis lines=middle,
          axis line style={->, thick, gray!80!black},
          xlabel=$x$,
          ylabel=$y$,
          xtick={1,3,\maxn},
          xticklabels={1,3,$n$},
          ytick={1},
          grid=none,
          samples=200,
          no markers,
          clip=false,
        ]

        % Draw the function y=x/sqrt(x^2-1)
        \addplot[thick, smooth, blue!85!black,domain=1.16:\maxx,samples=200] {gfunc(x)} node[above, pos=0.5, yshift=5mm, black] {$y=\frac{x}{\sqrt{x^2-1}}$};
        \addplot[draw=none, domain=3:\maxn, name path=A] {gfunc(x)} ;
        \addplot[draw=none, domain=3:\maxn, name path=line]{0};
        % Draw the vertical line at x=1
        \draw[densely dotted, thick, gray!70!black] (axis cs:1, 0) -- (axis cs:1, 2);
        % Draw the horizontal line at y=1
        \draw[densely dotted, thick, gray!70!black] (axis cs:0, 1) -- (axis cs:\maxn, 1);
        % Add a label for the origin
        \node[below left] at (axis cs:0, 0) {$O$};
        \draw[densely dotted, thick, gray!70!black] (axis cs:\maxn, 0) -- (axis cs:\maxn, 1);

        \addplot[pattern=north west lines, pattern color=gray!40!white] fill between [of=A and line];

        \draw [pattern=north east lines, pattern color=blue!85!black] (2-1,0) rectangle (2,{gfunc(2)});
        \draw [pattern=north east lines, pattern color=blue!85!black] (3-1,0) rectangle (3,{gfunc(3)});
        \draw [pattern=north east lines, pattern color=blue!85!black] (4-1,0) rectangle (4,{gfunc(4)});

        \draw [pattern=north east lines, pattern color=blue!85!black] (\maxn-3,0) rectangle (\maxn-2,{gfunc(\maxn-2)});
        \draw [pattern=north east lines, pattern color=blue!85!black] (\maxn-2,0) rectangle (\maxn-1,{gfunc(\maxn-1)});
        \draw [pattern=north east lines, pattern color=blue!85!black] (\maxn-1,0) rectangle   (\maxn,{gfunc(\maxn)});

      \end{axis}
    \end{tikzpicture}
    \caption{$f_n$の上からの評価．斜線部の面積が$f_n$を表し，グレーアウトの部分が$S_n$を表す．}
    \label{2001-1:fig:3}
  \end{figure}

  従って，まず\cref{2001-1:fig:3}から，$f(n)$を上から評価すると
  \begin{align}
    f(n) < S_n + g(2) + g(3) \label{eq:3}
  \end{align}
  である．

  \begin{figure}[H]
    \centering
    \pgfmathsetmacro\maxx{12}
    \pgfmathsetmacro\maxn{10}
    \begin{tikzpicture}
      \begin{axis}[
          xmin=0, xmax=\maxx,
          ymin=0, ymax=2.0,
          axis lines=middle,
          axis line style={->, thick, gray!80!black},
          xlabel=$x$,
          ylabel=$y$,
          xtick={1,2,3,\maxn},
          xticklabels={1,2,3,$n$},
          ytick={1},
          grid=none,
          samples=200,
          no markers,
          clip=false,
        ]

        % Draw the function y=x/sqrt(x^2-1)
        \addplot[thick, smooth, blue!85!black,domain=1.16:\maxx,samples=200] {gfunc(x)} node[above, pos=0.5, yshift=5mm, black] {$y=\frac{x}{\sqrt{x^2-1}}$};
        \addplot[draw=none, domain=3:\maxn, name path=A] {gfunc(x)} ;
        \addplot[draw=none, domain=3:\maxn, name path=line]{0};
        % Draw the vertical line at x=1
        \draw[densely dotted, thick, gray!70!black] (axis cs:1, 0) -- (axis cs:1, 2);
        % Draw the horizontal line at y=1
        \draw[densely dotted, thick, gray!70!black] (axis cs:0, 1) -- (axis cs:\maxn, 1);
        % Add a label for the origin
        \node[below left] at (axis cs:0, 0) {$O$};
        \draw[densely dotted, thick, gray!70!black] (axis cs:\maxn, 0) -- (axis cs:\maxn, 1);

        \addplot[pattern=north west lines, pattern color=gray!40!white] fill between [of=A and line];

        \draw [pattern=north east lines, pattern color=blue!85!black] (2,0) rectangle (2+1,{gfunc(2)});
        \draw [fill=gray!40] (2,{gfunc(2)}) -- (3,{gfunc(2)}) -- (3,{gfunc(3)}) -- cycle;
        \draw [pattern=north east lines, pattern color=blue!85!black] (3,0) rectangle (3+1,{gfunc(3)});
        \draw [fill=gray!40] (3,{gfunc(3)}) -- (4,{gfunc(3)}) -- (4,{gfunc(4)}) -- cycle;
        \draw [pattern=north east lines, pattern color=blue!85!black] (4,0) rectangle (4+1,{gfunc(4)});
        \draw [fill=gray!40] (4,{gfunc(4)}) -- (5,{gfunc(4)}) -- (5,{gfunc(5)}) -- cycle;

        \draw [pattern=north east lines, pattern color=blue!85!black] (\maxn-2,0) rectangle (\maxn-1,{gfunc(\maxn-2)});
        \draw [pattern=north east lines, pattern color=blue!85!black] (\maxn-1,0) rectangle (\maxn,{gfunc(\maxn-1)});
        \draw [pattern=north east lines, pattern color=blue!85!black] (\maxn,0)   rectangle (\maxn+1,{gfunc(\maxn)});

      \end{axis}
    \end{tikzpicture}
    \caption{$f_n$の下からの評価．斜線部の面積が$f_n$を表し，グレーアウトの部分が$S_n$を表す．厳しく評価をするため，$f_n$の上の三角形部分も評価に含める．}
    \label{2001-1:fig:4}
  \end{figure}

  次に\cref{2001-1:fig:4}で$y \ge g(x)$の三角形部分の面積が
  \begin{align}
    & \dfrac{g(3)-g(4)}{2} + \dfrac{g(4)-g(5)}{2} + \cdots + \dfrac{g(n-1)-g(n)}{2} \\
    & = \dfrac{g(3)-g(n)}{2}
  \end{align}
  だから，
  \begin{align}
    & S_n + g(2) + g(n) + \dfrac{g(3)-g(n)}{2} < f(n)                           \\
    & S_n + g(2) + \dfrac{g(3)}{2} + \dfrac{g(n)}{2} < f(n)  \label{eq:4}
  \end{align}
  である．

  以上\cref{eq:3,eq:4}をまとめて
  \begin{align}
    S_n + g(2) + \dfrac{g(3)}{2} + \dfrac{g(n)}{2} < f(n) < S_n + g(2) + g(3)
  \end{align}
  である．よって$A(n) = n-f(n)$より
  \begin{align}
    n-S_n - \left(g(2) + g(3)\right) < A(n) < n - S_n - \dfrac{g(n)}{2} - \left(g(2) + \dfrac{g(3)}{2}\right)
  \end{align}
  である．以下これを評価する．
  まず$S_n$は
  \begin{align}
    S_n = [\sqrt{x^2-1}]_3^{n} = \sqrt{n^2-1} - 2\sqrt{2}
  \end{align}
  だから，
  \begin{align}
    n-\sqrt{n^2-1} + \left(2\sqrt{2}-g(2)-g(3)\right) < A(n) < n - \sqrt{n^2-1} - \dfrac{n}{2\sqrt{n^2-1}} + \left(2\sqrt{2}-g(2)-\dfrac{g(3)}{2}\right) \label{eq:5}
  \end{align}
  となる．

  ここで新しく
  \begin{align}
    h(x) = x -\sqrt{x^2-1}
  \end{align}
  とおくと，
  \begin{align}
    h(x) = \dfrac{1}{x+\sqrt{x^2-1}}
  \end{align}
  より$2 \le x$では$h(x)$は単調減少である．
  そこで$n\to\infty$での\cref{eq:5}の両辺の振る舞いを考えると，
  \begin{align}
    \lim_{n\to\infty}\left[n-\sqrt{n^2-1}\right]
    & = \lim_{n\to\infty} \dfrac{1}{n+\sqrt{n^2-1}} \\
    & = 0                                          \\
    \lim_{n\to\infty}\left[n - \sqrt{n^2-1} - \dfrac{n}{2\sqrt{n^2-1}}\right]
    & = \dfrac{-1}{2}
  \end{align}
  だから，

  \begin{align}
    \left(2\sqrt{2}-g(2)-g(3)\right) < \lim_{n\to\infty}A(n) < \dfrac{-1}{2}+ \left(2\sqrt{2}-g(2)-\dfrac{1}{2}g(3)\right)
  \end{align}
  である．ここに
  \begin{align}
    g(2) & = \dfrac{2\sqrt{3}}{3} \\
    g(3) & = \dfrac{3\sqrt{2}}{4}
  \end{align}
  を代入して
  \begin{align}
    \dfrac{5\sqrt{2}}{4} - \dfrac{2\sqrt{3}}{3} < \lim_{n\to\infty}A(n) < \dfrac{-1}{2} + \dfrac{13\sqrt{2}}{8} - \dfrac{2\sqrt{3}}{3}
  \end{align}
  である．
  \begin{align}
    \sqrt{2} & = 1.414... \\
    \sqrt{3} & = 1.732...
  \end{align}
  だから，
  \begin{align}
    \dfrac{5\sqrt{2}}{4} - \dfrac{2\sqrt{3}}{3}
    & >  \dfrac{5\cdot 1.414}{4} - \dfrac{2\cdot 1.733}{3} \\
    & = 0.612...
  \end{align}
  および
  \begin{align}
    \dfrac{-1}{2} + \dfrac{13\sqrt{2}}{8} - \dfrac{2\sqrt{3}}{3}
    & < \dfrac{-1}{2} + \dfrac{13\cdot 1.415}{8} - \dfrac{2\cdot 1.732}{3} \\
    & = 0.644...
  \end{align}
  だから，
  \begin{align}
    0.612 < \lim_{n\to\infty}A(n) < 0.645
  \end{align}
  である．

  従って，任意の自然数$n$で$A_n \ge i/10$を満たすためには，
  \begin{align}
    i\le 6
  \end{align}
  が必要．

  また，$h(x)$の単調減少性から，
  \begin{align}
    A(n)
    & > h(n) + \left(2\sqrt{2}-g(2)-g(3)\right)                  \\
    & > \lim_{n\to\infty}h(n) + \left(2\sqrt{2}-g(2)-g(3)\right) \\
    & = \left(2\sqrt{2}-g(2)-g(3)\right)                         \\
    & > 0.612...
  \end{align}
  だから，$i\le 6$であれば常に$A(n) \ge i/10$を満たし十分である．
  以上から求める最大の$i$は
  \begin{align}
    i=6
  \end{align}
  である．$\cdots$(答)

  \vspace{10pt}
  {\bf [解説]}

  \newpage
\end{multicols}
\end{document}
